
\section{امنیت شبکه به زبان ساده (\lr{Network Security})}
\subsection{امنیت شبکه یعنی چی؟}
خیلی ساده بگیم: هدف امنیت شبکه اینه که داده‌ها \textbf{محرمانه} باشن (هر کسی نبینه)، \textbf{سالم} بمونن (هیچ‌کس وسط راه تغییرشون نده)، و \textbf{دردسترس} باشن (وقتی لازم داریم قطع نباشن). برای رسیدن به این سه‌تا هدف، یه سری سیاست و روش و ابزار طراحی می‌کنیم که جلوی تهدیدها رو بگیرن.

نکته تاریخی: اوایل اینترنت، اکثر پروتکل‌ها با فرض «اعتماد بین کاربرا» ساخته شدن و خیلی از ویژگی‌های امنیتی رو نداشتن؛ بعدها که استفاده عمومی شد، نیاز به راهکارهای دفاعی جدی‌تر احساس شد.

\subsection{بدافزارها (\lr{Malware}) با مثال‌های دم‌دستی}
\begin{itemize}
  \item \textbf{ویروس (\lr{Virus}):} یه برنامه مزاحم که خودش رو به یه برنامه اجرایی دیگه {می‌چسبونه} و برای پخش شدن لازم داره اون برنامه اجرا بشه.
  \item \textbf{کِرم (\lr{Worm}):} مستقل از بقیه‌ست، خودش می‌تونه {خودکار} از یه کامپیوتر به کامپیوتر دیگه تکثیر بشه؛ معمولاً از نقطه‌ضعف‌های شبکه بدون دخالت انسان استفاده می‌کنه.
  \item \textbf{\lr{Spyware}:} نرم‌افزارهای فضول که فعالیت‌های کاربر (مثل کلیدهایی که می‌زنه) رو ضبط می‌کنن و می‌فرستن برای مهاجم.
  \item \textbf{\lr{Botnet}:} وقتی کلی دستگاه آلوده زیر نظر مهاجم قرار بگیرن و به شکل هماهنگ برای حمله استفاده بشن. هر دستگاه آلوده می‌شه یه \lr{bot}.
\end{itemize}

\subsection{حملات رایج با توضیح کوتاه}
\paragraph{حمله منع سرویس \lr{Denial of Service (DoS)}}
هدف اینه که سرویس (مثلاً یه وب‌سایت) برای کاربرای مجاز \textbf{از کار بیفته}. چطوری؟ با فرستادن حجم عظیمی از درخواست‌های تقلبی تا \lr{CPU} یا پهنای‌باند سرور اشباع بشه و نتونه به درخواست‌های واقعی جواب بده.

اگه همین کار از \textbf{چندین} نقطه در دنیا و معمولاً از طریق یه \lr{Botnet} انجام بشه، می‌گیم \lr{DDoS} که مقابله باهاش سخت‌تره.

\paragraph{شنود بسته‌ها \lr{Packet Sniffing}}
توی بعضی شبکه‌ها (خصوصاً بی‌سیم) اگه کارت شبکه رو بذاریم روی حالت \lr{promiscuous mode}، می‌تونه \textbf{همه ترافیک عبوری} رو ببینه، نه فقط ترافیک خودش. مهاجم می‌تونه از داخل بسته‌ها اطلاعات حساسی مثل نام کاربری و گذرواژه رمز نشده رو \textbf{استخراج} کنه.

\paragraph{جعل هویت \lr{IP Spoofing}}
مهاجم آدرس \lr{IP} مبدأ رو توی بسته‌ها \textbf{دستکاری} می‌کنه تا یا منبعِ قابل‌اعتماد جا بزنه و دسترسی خاص بگیره، یا هویت واقعی خودش رو پنهان کنه.

\section{چرا شبکه رو لایه‌لایه می‌کنیم؟}
سیستم‌های شبکه‌ای \textbf{پیچیده} هستن. برای مدیریت این پیچیدگی، کارها رو به لایه‌های مجزا تقسیم می‌کنیم. هر لایه وظیفه مشخصی داره، از خدمات لایه پایینی استفاده می‌کنه و به لایه بالایی سرویس می‌ده. دو مزیت مهم:
\begin{itemize}
  \item \textbf{ساختار واضح:} کار هر لایه معلومه و مدیریت راحت‌تر می‌شه.
  \item \textbf{ماژولار بودن:} می‌تونیم یه لایه رو عوض یا به‌روزرسانی کنیم \textbf{بدون} این‌که بقیه لایه‌ها رو دست بزنیم. مثلاً می‌شه تکنولوژی \lr{WiFi} رو در لایه پیوند داده عوض کرد بدون این‌که لازم باشه مرورگر وب (لایه کاربرد) تغییر کنه.
\end{itemize}

\section{پشته پروتکل اینترنت (\lr{Internet Protocol Stack}) در عمل}
در عمل، اینترنت معمولاً با یه مدل \textbf{۵ لایه} دیده می‌شه. وقتی داده از بالا به پایین می‌ره، هر لایه اطلاعات خودش رو \textbf{کپسوله} (\lr{Encapsulation}) می‌کنه و به داده اضافه می‌شه.

\subsection{لایه کاربرد (\lr{Application Layer})}
\textbf{تعامل مستقیم با کاربر} از این‌جا انجام می‌شه. برنامه‌های کاربردی شبکه از پروتکل‌هایی مثل \lr{HTTP} (برای وب)، \lr{SMTP} (برای ایمیل)، \lr{FTP} (برای انتقال فایل) استفاده می‌کنن.

\textbf{واحد داده:} \lr{Message} (پیام)

\subsection{لایه انتقال (\lr{Transport Layer})}
ارتباط منطقی بین فرآیندهای در حال اجرا روی میزبان مبدأ و مقصد رو برقرار می‌کنه. دو پروتکل اصلی:
\begin{itemize}
  \item \lr{TCP}: اتصال‌گرا و \textbf{قابل اعتماد}. مرتب‌سازی و تحویل \textbf{صحیح و کامل} داده رو تضمین می‌کنه؛ مناسب برای وب، ایمیل و هر جا که \textbf{صحت} مهم‌تر از سرعته.
  \item \lr{UDP}: بدون اتصال و \textbf{سریع‌تر} ولی بدون تضمین تحویل. مناسب برای استریم ویدیو و بازی آنلاین که \textbf{تأخیر کم} مهم‌تره.
\end{itemize}

\textbf{واحد داده:} \lr{Segment} (قطعه)

\subsection{لایه شبکه (\lr{Network Layer})}
\textbf{مسیر دادن} بسته‌ها از کامپیوتر مبدأ به مقصد در سراسر اینترنت با \lr{IP Addressing}. پروتکل‌های مثال: \lr{IP} و پروتکل‌های مسیریابی مثل \lr{OSPF}.

\textbf{واحد داده:} \lr{Datagram} (دیتاگرام)

\subsection{لایه پیوند داده (\lr{Link Layer})}
انتقال داده بین دو گره همسایه در یک شبکه محلی. آدرس‌دهی فیزیکی (\lr{MAC Address}). مثال پروتکل‌ها: \lr{Ethernet} برای شبکه‌های سیمی و \lr{WiFi} برای بی‌سیم.

\textbf{واحد داده:} \lr{Frame} (قاب)

\subsection{لایه فیزیکی (\lr{Physical Layer})}
انتقال \textbf{بیت‌های خام} (صفر و یک) روی رسانه فیزیکی مثل کابل مسی، فیبر نوری یا امواج رادیویی.

\section{چند مفهوم کلیدی}
\subsection{کپسوله‌سازی (\lr{Encapsulation})}
فرض کن می‌خوای هدیه بفرستی: اول خود هدیه (\lr{Message}) رو می‌ذاری توی یه جعبه (می‌شه \lr{Segment})، بعد جعبه رو می‌ذاری توی یه کارتن بزرگ‌تر با آدرس (\lr{Datagram})، آخر سر کارتن رو می‌دی به پست محلی که برچسب محلی می‌زنه (\lr{Frame}). هر مرحله یه «روکش» بهش اضافه می‌شه تا مسیرش رو پیدا کنه.

\subsection{نام واحدهای داده در هر لایه یک‌جا}
\begin{center}
  \begin{tabular}{|c|c|}
    \hline
    لایه & نام واحد داده \\
    \hline
    کاربرد (\lr{Application}) & \lr{Message} \\
    انتقال (\lr{Transport}) & \lr{Segment} \\
    شبکه (\lr{Network}) & \lr{Datagram} \\
    پیوند داده (\lr{Link}) & \lr{Frame} \\
    فیزیکی (\lr{Physical}) & بیت‌های خام \\
    \hline
  \end{tabular}
\end{center}

\section{تاریخچه خیلی خلاصه اینترنت}
\subsection{دهه ۱۹۶۰ تا ۱۹۷۰}
تحقیقات بنیادی روی شبکه‌های \lr{packet-switching} شروع شد و شبکه \lr{ARPAnet} {(وزارت دفاع آمریکا)} راه افتاد؛ بیشتر کاربرد پژوهشی و نظامی داشت.

\subsection{دهه ۱۹۷۰ تا ۱۹۸۰}
شبکه‌ها زیاد شدن و نیاز بود \textbf{شبکه‌های مختلف} به هم وصل بشن (\lr{Internetworking}). پروتکل‌های پایه مثل \lr{TCP/IP} طراحی شدن.

\subsection{دهه ۱۹۸۰ تا ۱۹۹۰}
\lr{TCP/IP} استاندارد رسمی شد. سرویس‌های پایه مثل ایمیل (\lr{SMTP})، انتقال فایل (\lr{FTP}) و سیستم نام دامنه (\lr{DNS}) شکل گرفتن.

\subsection{دهه ۱۹۹۰ تا ۲۰۰0}
با ظهور \lr{Web} و مرورگرهای گرافیکی، اینترنت از ابزار آکادمیک تبدیل شد به پدیده‌ای جهانی و تجاری.

\subsection{۲۰۰۵ تا امروز}
انفجار دستگاه‌های همراه، شبکه‌های اجتماعی و سرویس‌های \lr{Cloud}. اینترنت تقریباً به همه جنبه‌های زندگی وارد شده.

\section{مدل مرجع \lr{ISO/OSI} در یک نگاه}
مدل \lr{ISO/OSI} یه چارچوب \textbf{۷ لایه}‌ای و بیشتر \textbf{مفهومی}ه برای طراحی شبکه. نسبت به مدل ۵ لایه اینترنت، دو لایه اضافه داره: \lr{Presentation} (ارائه) و \lr{Session} (نشست). در عمل، خیلی وقت‌ها این دو لایه با لایه کاربرد \lr{(Application)} \textbf{ادغام} می‌شن و پیاده‌سازی واقعی بیشتر بر مبنای مدل \lr{TCP/IP} انجام می‌گیره.